\documentclass[11pt,letterpaper]{article}

\usepackage[margin=1in]{geometry}
\usepackage{amsmath,amssymb,amsthm}
\usepackage{mathtools}
\usepackage{booktabs}
\usepackage{xcolor}
\usepackage{hyperref}
\hypersetup{colorlinks=true,linkcolor=blue!50!black,urlcolor=blue!60!black,citecolor=blue!50!black}

\newtheorem{theorem}{Theorem}[section]
\newtheorem{lemma}[theorem]{Lemma}
\newtheorem{proposition}[theorem]{Proposition}
\newtheorem{corollary}[theorem]{Corollary}
\newtheorem{definition}[theorem]{Definition}
\newtheorem{remark}[theorem]{Remark}

\newcommand{\HH}{\mathbb{H}}
\newcommand{\ZZ}{\mathbb{Z}}
\newcommand{\QQ}{\mathbb{Q}}
\newcommand{\CC}{\mathbb{C}}
\newcommand{\SLZ}{\mathrm{SL}_2(\ZZ)}
\newcommand{\Dhat}{\widehat{\Delta}}
\newcommand{\eps}{\varepsilon}
\DeclareMathOperator{\im}{Im}

\title{A Single Linear Functional for Periods and Quasi-Periods\\[0.2em]
       of Weight-12 Modular Forms\\[0.4em]
       \large (the $\dim S_k = 1$ setting)}
\author{}
\date{\today}

\begin{document}
\maketitle

\begin{abstract}
For any $f \in S_{12}^!$ on $\SLZ$ (the weakly holomorphic forms of weight $12$ with vanishing constant Fourier coefficient; this includes the unique normalised holomorphic cusp form $\Delta \in S_{12} \subset S_{12}^!$) we exhibit an explicit linear functional $f \mapsto (\omega^+(f),\,\omega^-(f)) \in \CC^2$ given by integration of $f$ along two finite line segments in the upper half-plane against universal $\QQ$-polynomial kernels.  Specialising to $\Delta$ recovers Eichler--Shimura's classical periods $\omega^\pm(\Delta)$; specialising to the unique $\Dhat \in S_{12}^!$ with $\Dhat(\tau) = q^{-1} + O(q^2)$ produces Brown's quasi-periods $\eta^\pm(\Dhat)$.  Same formula, same kernels, different $f$.  The construction is the explicit Bernoulli-operator unpacking of the Diamantis--Rolen finite-endpoint cocycle; the functional is basepoint-independent and lockstep across nine equivalent extractions, matching Brown's reported values to relative error $< 6\times 10^{-32}$.  Treating $\Delta$ and $\Dhat$ on the same footing makes this the natural starting point for extending period-polynomial theory to fully meromorphic modular forms.
\end{abstract}

\section{Introduction and main result}\label{sec:intro}

The period polynomial of a weight-$k$ holomorphic cusp form $f \in S_k$ on $\Gamma = \SLZ$ is the polynomial
\[
  r_f(X) \;:=\; \int_0^{i\infty}\! f(\tau)\,(\tau - X)^{k-2}\,d\tau
              \;\in\; V_{k-2}\otimes\CC,
\]
which decomposes into even and odd parts $r_f^\pm$ whose coefficients are $\QQ$-rational multiples of $\omega^\pm(f)\,(2\pi i)^\bullet$, with the periods $\omega^\pm(f)$ in turn expressible as $\QQ$-multiples of critical $L$-values $L(f, s)$.  The Eichler--Shimura picture realises $\omega^\pm(f)$ as $\CC$-linear functionals of $f$ via a contour integral, and the resulting coordinates carry deep arithmetic.

For weakly holomorphic $f \in S_k^!$ with a cusp pole, the integral $\int_0^{i\infty}$ diverges and the classical construction fails.
Nevertheless Brown \cite{Brown2017} defines and numerically computes emph{quasi-periods} $\eta^\pm(f)$ for the unique form $\Dhat \in S_{12}^!$ satisfying $\Dhat(\tau) = q^{-1} + O(q^2)$ (equivalently $\Dhat = \Delta\cdot(J^2 + 24J - 393444)$ with $J = j - 744$ the normalised Hauptmodul).  The qualifier $+\,O(q^2)$ is essential for uniqueness: the $q^0$ and $q^1$ coefficients of $\Dhat$ both vanish, which together with the principal part $q^{-1}$ pins $\Dhat$ down within $S_{12}^!$ (any $\Dhat + c\,\Delta$ would otherwise share its principal part).

Brown's construction uses the Diamantis--Rolen finite-endpoint cocycle
\cite{DR2017}: replacing the divergent contour $0 \to i\infty$ with a
finite segment $\gamma^{-1}\tau_0 \to \tau_0$ in $\HH$ produces a genuine
$1$-cocycle $C^f$ of $\Gamma$ whose cohomology class encodes
$\omega^\pm(f)$ for holomorphic $f$ and $\eta^\pm(f)$ for weakly
holomorphic $f$, on the same footing.  Brown's $\omega^\pm$ and
$\eta^\pm$ thus come from a single procedure.  The present note unpacks
this procedure to the level of an explicit linear functional, fully
transparent at the coefficient level.

\paragraph{Conventions.}
Fix $k = 12$, $n = k - 2 = 10$.  Let $V_n = \bigoplus_{i+j=n}\QQ X^i Y^j$
with the right $\Gamma$-action $(X,Y)|_\gamma = (aX + bY, cX + dY)$ for
$\gamma = \bigl(\begin{smallmatrix}a&b\\c&d\end{smallmatrix}\bigr)$.
Write $S = \bigl(\begin{smallmatrix}0&-1\\1&0\end{smallmatrix}\bigr)$ and
$T = \bigl(\begin{smallmatrix}1&1\\0&1\end{smallmatrix}\bigr)$.  We
write $[P]_{X^{10-j}Y^j}$ for the coefficient of $X^{10-j}Y^j$ in $P \in V_{10}$,
and set $q = e^{2\pi i \tau}$.  We use the convention that $S_k^!$ denotes
the weakly holomorphic forms of weight $k$ with vanishing constant
Fourier coefficient, so $S_k \subset S_k^!$.

Zagier's space $W_{12} \subset V_{10}$ of admissible period polynomials
at weight $12$ is $3$-dimensional, with the standard basis
\begin{equation}\label{eq:Zagier-basis}
\begin{aligned}
  P_0(X,Y) &\;=\; X^{10} - Y^{10},\\
  P_1(X,Y) &\;=\; X^8 Y^2 - 3 X^6 Y^4 + 3 X^4 Y^6 - X^2 Y^8,\\
  P_2(X,Y) &\;=\; 4 X^9 Y - 25 X^7 Y^3 + 42 X^5 Y^5 - 25 X^3 Y^7 + 4 X Y^9,
\end{aligned}
\end{equation}
splitting as $W_{12} = W_{12}^+ \oplus W_{12}^-$ with
$W_{12}^+ = \langle P_0, P_1\rangle$ (the $2$-dimensional even-in-$X$
sector, containing the Eisenstein direction $P_0$) and
$W_{12}^- = \langle P_2\rangle$ (the $1$-dimensional odd-in-$X$ cuspidal
sector) \cite{KZ1984}.  Brown's basis at $S$ \cite[\S 8.2]{Brown2017} is
$P^+_S = -\tfrac{36}{691} P_0 + P_1$ and $P^-_S = P_2$.

\paragraph{The main result.}
For each interior monomial $j \in \{1, 2, \ldots, 9\}$ let $\eps_j$
denote the parity sign --- $-$ if $j$ is odd, $+$ if $j$ is even --- and
let $c_S^j := [P^{\eps_j}_S]_{X^{10-j}Y^j} \in \QQ$ denote the
coefficient of $X^{10-j}Y^j$ in Brown's basis vector of the appropriate
parity.

\begin{theorem}[Periods and quasi-periods as a linear functional]\label{thm:main}
There exist universal polynomials $k_S^j, k_T^j \in \QQ[\tau]$ of degree
$\le 10$ in $\tau$, indexed by $j \in \{1, \ldots, 9\}$ and given
explicitly in Table~\ref{tab:kernels} below, with the following
property.  For any $\tau_0 \in \HH$ and any $f \in S_{12}^!$,
\begin{equation}\label{eq:main}
  \boxed{\;\;
  \omega^{\eps_j}(f) \;=\; \frac{(2\pi i)^{11}}{c_S^j}
    \left[\,\int_{-1/\tau_0}^{\tau_0}\!\!f(\tau)\,k_S^{j}(\tau)\,d\tau
       \;+\; \int_{\tau_0 - 1}^{\tau_0}\!\!f(\tau)\,k_T^{j}(\tau)\,d\tau\,\right].
  \;\;}
\end{equation}
The value $\omega^{\eps_j}(f)$ is independent of the basepoint $\tau_0$
and of the choice of interior monomial $j$ (within its parity sector).
Both contour integrals are finite for every $f \in S_{12}^!$, with no
regularisation, analytic continuation, or limit required.
\end{theorem}

Theorem~\ref{thm:main} is the analog of the classical
$\omega^\pm \propto L(f, s)\,(2\pi i)^\bullet$ formula in finite-contour
DR coordinates, valid uniformly for holomorphic and weakly holomorphic
$f$.  We prove it in \S\S\ref{sec:DR-cocycle}--\ref{sec:cohomology}:
existence of the kernels via the Bernoulli operator
(\S\ref{sec:kernels}), basepoint-independence via an explicit coboundary
identity (\S\ref{sec:cohomology}), $j$-independence via the lockstep
structure of $W_{12}$.  Two specialisations of immediate interest:

\begin{corollary}[Classical periods $\omega^\pm(\Delta)$]\label{cor:Delta}
For $f = \Delta \in S_{12} \subset S_{12}^!$, the unique normalised
holomorphic cusp form, \eqref{eq:main} reproduces the Eichler--Shimura
periods.  At $\tau_0 = 0.3 + 1.2i$ all nine extractions agree (relative
residual $\sim 10^{-24}$) on
\[
  \omega^+(\Delta) = -68\,916\,772.80959519475431\ldots, \qquad
  \omega^-(\Delta) = -5\,585\,015.379310401866877\ldots\,i,
\]
matching Brown \cite[\S 8.3]{Brown2017} to all $20$ reported digits.
\end{corollary}

\begin{corollary}[Quasi-periods $\eta^\pm(\Dhat)$]\label{cor:Dhat}
For $f = \Dhat = \Delta\cdot(J^2 + 24J - 393444) \in S_{12}^!$,
\eqref{eq:main} reproduces Brown's quasi-periods.  At
$\tau_0 = 0.3 + 1.2i$, all nine extractions agree (relative residual
$\sim 10^{-17}$) on
\[
  \eta^+(\Dhat) = 127\,202\,100\,647.1770947773\ldots, \qquad
  \eta^-(\Dhat) = 10\,276\,732\,343.6491327508\ldots\,i,
\]
matching Brown \cite[\S 8.3]{Brown2017} to all $20$ reported digits.
\end{corollary}

The integrand and contour structure are \emph{identical} in
Corollaries~\ref{cor:Delta} and~\ref{cor:Dhat}; only the modular form
$f$ under the integral sign changes.  This is the precise sense in
which the DR construction places $\omega^\pm(\Delta)$ and
$\eta^\pm(\Dhat)$ on the same footing.

\begin{corollary}[Uniform applicability across $\CC\lbrack J\rbrack\cdot\Delta$]\label{cor:RJ}
For any $R \in \CC[J]$ with $a_{R\Delta}(0) = 0$, formula \eqref{eq:main}
computes $\omega^\pm(R\cdot\Delta)$ via the same universal kernels.  In
particular, the entire $\CC$-span of products $R(J)\,\Delta$ (with
vanishing constant term) maps linearly into
$H^1_{\rm cusp}(\Gamma; V_{10}\otimes\CC) = \CC[P^+]\oplus\CC[P^-]$,
landing on the pair $(\omega^+,\omega^-) \in \CC^2$ via \eqref{eq:main}.
\end{corollary}

\paragraph{Why this matters.}
A ``period'' is, in classical usage, the integral of an algebraic
differential along a cycle.  Eichler--Shimura realises
$\omega^\pm(\Delta)$ this way: as Mellin transforms of $\Delta$ on
$[0, i\infty]$, packaged into a polynomial in $V_{10}$.
Theorem~\ref{thm:main} preserves the period-as-integral picture but on
a finite contour in $\HH$, with universal polynomial kernel weights
$k_S^j, k_T^j$ in place of the bare $\tau^j$ moments.  For holomorphic
$\Delta$, taking $\tau_0 \to i\infty$ collapses the $T$-segment, sends
the $S$-segment to $[0, i\infty]$, and reduces the kernels to bare
monomials --- recovering classical Mellin
(Lemma~\ref{lem:Mellin-recovery}).  For weakly holomorphic $\Dhat$, the
limit is unavailable, but at every finite $\tau_0 \in \HH$ the formula
and its value are well-defined.  ``Quasi-periods'' are honestly periods
in the same geometric sense as the classical ones: integrals of a
modular form along explicit cycles in $\HH$, against explicit
polynomial kernels, with the contour bounded away from the cusp.  The
basepoint $\tau_0$ --- not analytic continuation, not a non-holomorphic
regulator --- is what tames the cusp pole.

\paragraph{Outline.}
\S\ref{sec:DR-cocycle} introduces the Diamantis--Rolen cocycle and its
moment expansion; \S\ref{sec:coboundary} introduces the coboundary
modification by a polynomial $P$ satisfying $P|_T - P = C^f_T$
(existence, uniqueness, Bernoulli closed form); \S\ref{sec:kernels}
establishes the universal kernels of Theorem~\ref{thm:main} and presents
Table~\ref{tab:kernels}; \S\ref{sec:specialisation} verifies the formula
numerically for $\Delta$ and $\Dhat$ and shows it produces an
exponentially convergent $q$-series for $\omega^\pm$;
\S\ref{sec:cohomology} interprets $\omega^\pm$ as $H^1_{\rm cusp}$
coordinates and proves the basepoint and $j$-independence;
\S\ref{sec:bfk} contrasts with BFK's regularisation; \S\ref{sec:outlook}
sketches the extension to fully meromorphic forms with interior poles
and offers a gauge-theoretic analogy as a thread for future work.

\section{The Diamantis--Rolen cocycle and its moment expansion}\label{sec:DR-cocycle}

\begin{definition}[Diamantis--Rolen finite-endpoint cocycle]\label{def:DR}
For $f \in S_{12}^!$ and $\tau_0 \in \HH$, define
$C^f \colon \Gamma \to V_{10}\otimes\CC$ by
\begin{equation}\label{eq:DR-def}
  C^f_\gamma(X, Y) \;:=\; (2\pi i)^{11}\!\int_{\gamma^{-1}\tau_0}^{\tau_0}\!
                          f(\tau)\,(X - \tau Y)^{10}\,d\tau,
\end{equation}
the contour being the straight line segment from $\gamma^{-1}\tau_0$ to
$\tau_0$ in $\HH$.
\end{definition}

\begin{lemma}\label{lem:DR-cocycle}
$C^f$ is a $1$-cocycle of $\Gamma$:
$C^f_{\gamma_1\gamma_2} = C^f_{\gamma_1}|_{\gamma_2} + C^f_{\gamma_2}$
for all $\gamma_1, \gamma_2 \in \Gamma$.
\end{lemma}

\begin{proof}
Path-additivity:
$\int_{(\gamma_1\gamma_2)^{-1}\tau_0}^{\tau_0} =
 \int_{(\gamma_1\gamma_2)^{-1}\tau_0}^{\gamma_2^{-1}\tau_0}
 + \int_{\gamma_2^{-1}\tau_0}^{\tau_0}$.
The substitution $\tau = \gamma_2\sigma$ in the first integral, combined
with weight-$12$ modularity of $f(\tau)\,d\tau$ (transforming with
weight $10$ under $\sigma \mapsto \gamma_2\sigma$, exactly cancelling
the weight-$-10$ slash on $(X-\tau Y)^{10}$), identifies it with
$C^f_{\gamma_1}|_{\gamma_2}(X, Y)/(2\pi i)^{11}$.
\end{proof}

A $1$-cocycle of $\Gamma$ is determined by its values at the generators
$S, T$.  Define the \emph{moment integrals}
\begin{equation}\label{eq:moments}
  N_j(f, \tau_0) := \int_{-1/\tau_0}^{\tau_0}\! f(\tau)\,\tau^j\,d\tau,
  \qquad
  M_j(f, \tau_0) := \int_{\tau_0 - 1}^{\tau_0}\! f(\tau)\,\tau^j\,d\tau,
  \qquad j = 0, 1, \ldots, 10.
\end{equation}
Expanding $(X - \tau Y)^{10}$ binomially in \eqref{eq:DR-def} gives
\begin{align}
  C_S^f(X, Y) &= (2\pi i)^{11}\sum_{j=0}^{10}(-1)^j\binom{10}{j}N_j\,X^{10-j}Y^j,
              \label{eq:CS-explicit}\\
  C_T^f(X, Y) &= (2\pi i)^{11}\sum_{j=0}^{10}(-1)^j\binom{10}{j}M_j\,X^{10-j}Y^j.
              \label{eq:CT-explicit}
\end{align}
For $f \in S_{12}^!$ the $q$-coefficients $a_f(n)$ grow at most as
$e^{4\pi\sqrt{n}}$ \cite{Brown2017}, while $|q|^n = e^{-2\pi n\im(\tau)}$
decays super-exponentially in $n$ for any $\tau \in \HH$.  The Fourier
series therefore converges absolutely on every compact subset of $\HH$,
and the moment integrals \eqref{eq:moments} are unambiguously finite for
any $\tau_0 \in \HH$.

\begin{remark}\label{rem:cuspidality-input}
The cuspidality input $a_f(0) = 0$ (built into the definition of
$S_{12}^!$) enters in exactly one place: it makes $M_0(f, \tau_0)
= \int_{\tau_0-1}^{\tau_0}\! f\,d\tau = a_f(0)$ (by periodicity of $f$),
so vanishing of $a_f(0)$ is equivalent to vanishing of the $X^{10}$
coefficient of $C^f_T$.  This is the solvability condition for the
coboundary equation in \S\ref{sec:coboundary}.
\end{remark}

\section{Cuspidality and the coboundary polynomial}\label{sec:coboundary}

The cocycle $C^f$ encodes more data than $\omega^\pm(f)$.  To extract
the periods we replace $C^f$ by a cohomologous \emph{cuspidal} cocycle
whose value at $T$ vanishes.  The price is a coboundary modification by
$P \in V_{10}\otimes\CC$ satisfying
\begin{equation}\label{eq:P-equation}
  P\big|_T \;-\; P \;=\; C^f_T. \tag{$\star$}
\end{equation}
The modified cocycle
\begin{equation}\label{eq:Cprime-def}
  C^{\prime f}_\gamma \;:=\; C^f_\gamma - (P|_\gamma - P)
\end{equation}
is cohomologous to $C^f$ and satisfies $C^{\prime f}_T = 0$.

\begin{proposition}[Existence and uniqueness of $P$]\label{prop:P-exists}
The operator $\delta_T \colon V_{10} \to V_{10}$, $P \mapsto P|_T - P$,
has $\ker(\delta_T) = \CC\cdot Y^{10}$ and
$\operatorname{im}(\delta_T) = \{Q \in V_{10} : [Q]_{X^{10}} = 0\}$.
Hence \eqref{eq:P-equation} is solvable iff $[C^f_T]_{X^{10}} = 0$,
which by Remark~\ref{rem:cuspidality-input} is equivalent to $a_f(0) =
0$ (automatic for $f \in S_{12}^!$).  The solution $P$ is unique modulo
$\CC\cdot Y^{10}$.
\end{proposition}

\begin{proof}
$T$ acts on $V_{10}$ as $X \mapsto X+Y$, $Y \mapsto Y$, fixing only
powers of $Y$, hence $\ker(\delta_T) = \CC\cdot Y^{10}$.  Writing
$P = \sum_{j=0}^{10}p_j X^{10-j}Y^j$ and expanding $(X+Y)^{10-j}$
binomially, the $X^{10-m}Y^m$ coefficient of $P|_T - P$ is
$\sum_{j=0}^{m-1}\binom{10-j}{m-j}p_j$.  This is upper-triangular in
$m$: the $m=0$ row is identically $0$ (requiring $[C^f_T]_{X^{10}}=0$),
and rows $m = 1, \ldots, 10$ determine $p_0, \ldots, p_9$ in turn, with
$p_{10}$ unconstrained.
\end{proof}

We fix the convention $p_{10} = 0$; alternative choices shift
$C^{\prime f}_S$ by $c(X^{10} - Y^{10})$ (since
$Y^{10}|_S - Y^{10} = X^{10} - Y^{10}$), a coboundary in $V_{10}$
leaving the cohomology class and the periods invariant.

\begin{proposition}[Bernoulli closed form for $P$]\label{prop:Bernoulli-P}
With $T = e^{Y\partial_X}$ acting on $V_{10}$ as a translation operator,
the unique solution of \eqref{eq:P-equation} with $p_{10} = 0$ is
\begin{equation}\label{eq:Bernoulli-P}
  P(X, Y) \;=\; \frac{1}{Y}\!\int_0^X\! C^f_T(t, Y)\,dt
              \;+\; \sum_{k=1}^{10}\frac{B_k}{k!}\,Y^{k-1}\,
                    \frac{\partial^{k-1} C^f_T}{\partial X^{k-1}}(X, Y),
\end{equation}
where $B_k$ are Bernoulli numbers ($B_0 = 1$, $B_1 = -\tfrac12$,
$B_2 = \tfrac16$, $B_4 = -\tfrac{1}{30}$, \ldots; $B_{2\ell+1} = 0$ for
$\ell \ge 1$).  The series terminates at $k = 10$, with at most six
nonzero Bernoulli terms.
\end{proposition}

\begin{proof}
Formally $\delta_T = e^{Y\partial_X} - 1$ as an operator, so
$\delta_T^{-1} = \sum_{k\ge 0}B_k(Y\partial_X)^{k-1}/k!$ by the
Bernoulli generating identity $z/(e^z - 1) = \sum_k B_k z^k/k!$.  The
$k=0$ summand is $(Y\partial_X)^{-1} = \partial_X^{-1}/Y$ (the
antiderivative in $X$ divided by $Y$), giving the integral term; the
$k \ge 1$ summands give the differential terms.  On a polynomial of
finite $X$-degree the operator series terminates.  Agreement of
\eqref{eq:Bernoulli-P} with the upper-triangular recursion in the proof
of Proposition~\ref{prop:P-exists} is verified to relative residual
$\sim 10^{-42}$ at $45$-digit working precision in \cite{VerifyBernoulli}.
\end{proof}

Substituting \eqref{eq:CT-explicit} into \eqref{eq:Bernoulli-P}
expresses each $p_l$ as an explicit $\QQ$-linear combination of the
moments $M_1, \ldots, M_{10}$:
\begin{equation}\label{eq:p-explicit}
  p_0 = -(2\pi i)^{11}\,M_1, \quad
  p_1 = 5\,(2\pi i)^{11}(M_1 + M_2), \quad
  p_2 = -(2\pi i)^{11}\!\left(\tfrac{15}{2}M_1 + \tfrac{45}{2}M_2 + 15 M_3\right),
\end{equation}
with $p_3, \ldots, p_9$ following the same pattern.

\begin{definition}[Modified cocycle value $C^{\prime f}_S$]\label{def:Cprime}
With $P$ as in Proposition~\ref{prop:Bernoulli-P} and using
$P|_S(X, Y) = P(Y, -X)$:
\begin{equation*}
  C^{\prime f}_S(X, Y) \;:=\; C^f_S(X, Y) \;-\; \bigl[\,P(Y, -X) - P(X, Y)\,\bigr].
\end{equation*}
This is a polynomial in $V_{10}\otimes\CC$ whose coefficients are
explicit linear combinations of the moments $\{M_j, N_j\}_{j=0}^{10}$ of
$f$, with rational/Bernoulli weights and overall factor $(2\pi i)^{11}$.
\end{definition}

By construction $C^{\prime f}_T = 0$, so $C^{\prime f}_S$ is the
$S$-value of a cuspidal $1$-cocycle and admits Brown's decomposition
\begin{equation}\label{eq:Brown-decomp}
  C^{\prime f}_S \;=\; \omega^+(f)\,P^+_S \;+\; \omega^-(f)\,P^-_S \;+\; c(f)\,(X^{10} - Y^{10}),
\end{equation}
with the $c(f)$ ambiguity absorbing the $p_{10} = 0$ convention.

\section{The universal polynomial kernels}\label{sec:kernels}

Definitions~\ref{def:DR},~\ref{def:Cprime} and
Proposition~\ref{prop:Bernoulli-P} express $C^{\prime f}_S$ as an
explicit $\QQ$-linear combination of the moments $\{N_j, M_j\}_{j=0}^{10}$
of $f$.  Reading off interior-monomial coefficients
$[C^{\prime f}_S]_{X^{10-j}Y^j}$ from \eqref{eq:Brown-decomp} and dividing
by $c_S^j$ gives $\omega^{\eps_j}(f)$ as a $\QQ$-linear functional of the
moments.  Each $N_j$ (resp.\ $M_j$) is an integral of $f$ against $\tau^j$
on the $S$-segment (resp.\ $T$-segment), so chaining gives an integral of
$f$ against an explicit polynomial in $\tau$.

\begin{theorem}[Universal kernel polynomials]\label{thm:kernels}
For each interior $j \in \{1, 2, \ldots, 9\}$, define
\begin{align}
  k_S^{j}(\tau) &\;:=\; (-1)^j\binom{10}{j}\,\tau^j, \label{eq:kS-def}\\
  k_T^{j}(\tau) &\;:=\; \sum_{r=1}^{10}\bigl[\,\beta_{j,r} \;-\; (-1)^j\,\beta_{10-j,r}\,\bigr]\,\tau^r, \label{eq:kT-def}
\end{align}
where $\beta_{l, r} \in \QQ$ are the Bernoulli-operator weights from
\eqref{eq:Bernoulli-P} computing the $X^{10-l}Y^l$ coefficient of $P$
from the moment $M_r$,
\begin{equation}\label{eq:beta-def}
  \beta_{l, r} \;=\;
    \frac{(-1)^r\,\binom{10}{r}\,B_{l-r+1}\,(10-r)!}{(l-r+1)!\,(10-l)!}
  \quad\text{for } 0 \le l \le 9,\; 1 \le r \le l+1,
\end{equation}
and $\beta_{l, r} = 0$ otherwise (in particular $\beta_{10, r} = 0$ for
all $r$, encoding $p_{10} = 0$).  With these kernels, identity
\eqref{eq:main} of Theorem~\ref{thm:main} holds for each interior $j$.
\end{theorem}

\begin{proof}
By Definition~\ref{def:Cprime},
$[C^{\prime f}_S]_{X^{10-j}Y^j} = [C^f_S]_{X^{10-j}Y^j}
 - [P(Y, -X) - P(X, Y)]_{X^{10-j}Y^j}$.
The first term is $(2\pi i)^{11}(-1)^j\binom{10}{j}\,N_j$ by
\eqref{eq:CS-explicit}; since
$N_j = \int_{-1/\tau_0}^{\tau_0}\!f(\tau)\tau^j\,d\tau$, the $S$-segment
contribution to
$\omega^{\eps_j}(f) = [C^{\prime f}_S]_{X^{10-j}Y^j}/c_S^j$ is
$(2\pi i)^{11}/c_S^j \cdot \int_{-1/\tau_0}^{\tau_0}\!f(\tau)\,k_S^j(\tau)\,d\tau$
with $k_S^j$ as in \eqref{eq:kS-def}.

For the second term, write $P = \sum_{l}p_l X^{10-l}Y^l$ and
$P(Y, -X) = \sum_l p_l(-1)^l X^l Y^{10-l}$.  The $X^{10-j}Y^j$
coefficient of $P(Y, -X)$ is $(-1)^{10-j}p_{10-j} = (-1)^j p_{10-j}$
(since $(-1)^{10} = 1$), so $[P(Y, -X) - P(X, Y)]_{X^{10-j}Y^j} =
(-1)^j p_{10-j} - p_j$.  Each $p_l$ is a $\QQ$-linear combination of
moments $p_l = (2\pi i)^{11}\sum_r \beta_{l,r}\,M_r$ via
\eqref{eq:Bernoulli-P}; direct expansion gives the formula
\eqref{eq:beta-def}.  Therefore
\[
  [P(Y, -X) - P(X, Y)]_{X^{10-j}Y^j} \;=\;
    (2\pi i)^{11}\sum_{r=1}^{10}\bigl[\,(-1)^j\beta_{10-j, r} - \beta_{j, r}\,\bigr]\,M_r,
\]
and the $T$-segment contribution to $\omega^{\eps_j}(f)$ is
$-(2\pi i)^{11}/c_S^j\cdot\int_{\tau_0-1}^{\tau_0}\!f(\tau)\,Q(\tau)\,d\tau$
where $Q(\tau) = \sum_r [(-1)^j\beta_{10-j, r} - \beta_{j, r}]\tau^r$.
The negative sign flips this into $k_T^j(\tau)$ as in \eqref{eq:kT-def}.
Adding the $S$- and $T$-contributions gives \eqref{eq:main}.
Independence of $\tau_0$ and $j$ is Theorem~\ref{thm:cohomology} below.
\end{proof}

The kernels are universal: they depend only on $j$, not on $f$ or
$\tau_0$.  Computing them via sympy and the Bernoulli formula
\eqref{eq:Bernoulli-P}--\eqref{eq:beta-def} \cite{KernelScript} gives
Table~\ref{tab:kernels}.

\begin{table}[h]\centering\renewcommand{\arraystretch}{1.25}
\begin{tabular}{cccll}
\toprule
$j$ & $\eps_j$ & $c_S^j$ & $k_S^{j}(\tau)$ & $k_T^{j}(\tau)$ \\
\midrule
$1$ & $-$ & $4$    & $-10\,\tau$        & $-5\,\tau + \tfrac{7}{2}\tau^{2} + 5\,\tau^{4} - 7\,\tau^{6} + \tfrac{15}{2}\tau^{8} - 5\,\tau^{9} + \tau^{10}$ \\
$2$ & $+$ & $1$    & $45\,\tau^{2}$     & $-9\,\tau + \tfrac{45}{2}\tau^{2} - 5\,\tau^{3} - 21\,\tau^{5} + 30\,\tau^{7} - \tfrac{45}{2}\tau^{8} + 5\,\tau^{9}$ \\
$3$ & $-$ & $-25$  & $-120\,\tau^{3}$   & $40\,\tau^{2} - 60\,\tau^{3} - 5\,\tau^{4} + 70\,\tau^{6} - 60\,\tau^{7} + 15\,\tau^{8}$ \\
$4$ & $+$ & $-3$   & $210\,\tau^{4}$    & $12\,\tau - 105\,\tau^{3} + 105\,\tau^{4} + 63\,\tau^{5} - 105\,\tau^{6} + 30\,\tau^{7}$ \\
$5$ & $-$ & $42$   & $-252\,\tau^{5}$   & $-42\,\tau^{2} + 210\,\tau^{4} - 252\,\tau^{5} + 84\,\tau^{6}$ \\
$6$ & $+$ & $3$    & $210\,\tau^{6}$    & $-12\,\tau + 105\,\tau^{3} - 105\,\tau^{4} - 63\,\tau^{5} + 105\,\tau^{6} - 30\,\tau^{7}$ \\
$7$ & $-$ & $-25$  & $-120\,\tau^{7}$   & $40\,\tau^{2} - 60\,\tau^{3} - 5\,\tau^{4} + 70\,\tau^{6} - 60\,\tau^{7} + 15\,\tau^{8}$ \\
$8$ & $+$ & $-1$   & $45\,\tau^{8}$     & $9\,\tau - \tfrac{45}{2}\tau^{2} + 5\,\tau^{3} + 21\,\tau^{5} - 30\,\tau^{7} + \tfrac{45}{2}\tau^{8} - 5\,\tau^{9}$ \\
$9$ & $-$ & $4$    & $-10\,\tau^{9}$    & $-5\,\tau + \tfrac{7}{2}\tau^{2} + 5\,\tau^{4} - 7\,\tau^{6} + \tfrac{15}{2}\tau^{8} - 5\,\tau^{9} + \tau^{10}$ \\
\bottomrule
\end{tabular}
\caption{The nine universal polynomial-kernel pairs $(k_S^j, k_T^j)$ of
Theorem~\ref{thm:kernels}, with Brown's basis values
$c_S^j = [P^{\eps_j}_S]_{X^{10-j}Y^j}$.  Computed exactly via
$\delta_T^{-1}$ in sympy rationals \cite{KernelScript}.  Pairs $(1, 9)$
and $(3, 7)$ share $k_T$ identically; pairs $(2, 8)$ and $(4, 6)$ share
$k_T$ up to overall sign (matching the sign flip in $c_S^j$); $j = 5$
is self-dual.}
\label{tab:kernels}
\end{table}

\begin{remark}\label{rem:S-kernel-monomial}
The $S$-kernel $k_S^j(\tau)$ is a pure monomial of degree $j$ --- it
comes directly from the bare $\tau^j$ moment $N_j$.  All complexity sits
in the $T$-kernel $k_T^j(\tau)$, which encodes the Bernoulli-rational
coboundary correction.
\end{remark}

\section{Specialisation to \texorpdfstring{$\Delta$ and $\Dhat$}{Delta and Dhat}}\label{sec:specialisation}

We now specialise Theorem~\ref{thm:main} to $f = \Delta$ and $f = \Dhat$
and verify Corollaries~\ref{cor:Delta},~\ref{cor:Dhat} numerically.

\begin{corollary}[Explicit formula for $\omega^\pm(\Delta)$]\label{cor:Delta-explicit}
For $\Delta \in S_{12}$, any $\tau_0 \in \HH$, and any interior $j \in \{1, \ldots, 9\}$,
\begin{equation}\label{eq:Delta-main}
  \omega^{\eps_j}(\Delta) \;=\; \frac{(2\pi i)^{11}}{c_S^j}
    \left[\,\int_{-1/\tau_0}^{\tau_0}\!\!\Delta(\tau)\,k_S^j(\tau)\,d\tau
        + \int_{\tau_0 - 1}^{\tau_0}\!\!\Delta(\tau)\,k_T^j(\tau)\,d\tau\,\right].
\end{equation}
\end{corollary}

\begin{corollary}[Explicit formula for $\eta^\pm(\Dhat)$]\label{cor:Dhat-explicit}
For $\Dhat \in S_{12}^!$, any $\tau_0 \in \HH$, and any interior $j \in \{1, \ldots, 9\}$,
\begin{equation}\label{eq:Dhat-main}
  \eta^{\eps_j}(\Dhat) \;=\; \frac{(2\pi i)^{11}}{c_S^j}
    \left[\,\int_{-1/\tau_0}^{\tau_0}\!\!\Dhat(\tau)\,k_S^j(\tau)\,d\tau
        + \int_{\tau_0 - 1}^{\tau_0}\!\!\Dhat(\tau)\,k_T^j(\tau)\,d\tau\,\right].
\end{equation}
\end{corollary}

The kernels in \eqref{eq:Delta-main} and \eqref{eq:Dhat-main} are
\emph{identical}; only the modular form $f$ under the integral changes.

\paragraph{Numerical verification.}
At $\tau_0 = 0.3 + 1.2i$, the integrals in \eqref{eq:Delta-main} and
\eqref{eq:Dhat-main} are computed at $45$-digit \texttt{mpmath} precision
in \cite{Notebook}.  Defining the DR-regularised $L$-value
$L_{\rm reg}(f, j+1)$ by $\omega^{\eps_j}(f) c_S^j /
[(2\pi)^{10-j}(-i)^j\binom{10}{j}\,j!]$ (which agrees with the classical
$L(\Delta, j+1)$ for $f = \Delta$, by Lemma~\ref{lem:Mellin-recovery}
below):

\begin{center}
\begin{tabular}{cccl}
\toprule
$j$ & $\eps_j$ & $L_{\rm reg}(\Delta, j+1)$ & extracted $\omega^{\eps_j}(\Delta)$ \\
\midrule
1  & $-$ & $0.14637454209126598\ldots$  & $-5\,585\,015.379310401\ldots\,i$ \\
2  & $+$ & $0.31524156588099308\ldots$  & $-68\,916\,772.80959519\ldots$ \\
3  & $-$ & $0.50161764751090221\ldots$  & $-5\,585\,015.379310401\ldots\,i$ \\
4  & $+$ & $0.66670918843400364\ldots$  & $-68\,916\,772.80959519\ldots$ \\
5  & $-$ & $0.79212283864603056\ldots$  & $-5\,585\,015.379310401\ldots\,i$ \\
6  & $+$ & $0.87735412538866091\ldots$  & $-68\,916\,772.80959519\ldots$ \\
7  & $-$ & $0.93070703029812609\ldots$  & $-5\,585\,015.379310401\ldots\,i$ \\
8  & $+$ & $0.96212645969442586\ldots$  & $-68\,916\,772.80959519\ldots$ \\
9  & $-$ & $0.97980908825122051\ldots$  & $-5\,585\,015.379310401\ldots\,i$ \\
\bottomrule
\end{tabular}
\end{center}

All nine rows lockstep to relative residual $< 6\times 10^{-42}$
\cite{VerifyL}.  The same nine integrals with $\Dhat$ in place of
$\Delta$ produce $\eta^+, \eta^-$ in lockstep to relative residual
$\sim 10^{-17}$ \cite{NumericScript}, agreeing with Brown's reported
values to all $20$ reported digits (Corollary~\ref{cor:Dhat}); the
slightly reduced precision reflects the larger absolute magnitudes
($|\eta^\pm| \sim 10^{11}$ versus $|\omega^\pm| \sim 10^7$).

\begin{lemma}[Mellin-transform recovery for $\Delta$]\label{lem:Mellin-recovery}
As $\tau_0 \to i\infty$, the $T$-segment integral in \eqref{eq:Delta-main}
vanishes (since all $M_j(\Delta, \tau_0) \to 0$ by cuspidality of
$\Delta$), the $S$-segment becomes $[0, i\infty]$, and the Mellin transform
$\int_0^{i\infty}\!\Delta(\tau)\tau^j\,d\tau = i^{j+1}j!/(2\pi)^{j+1}\cdot L(\Delta, j+1)$
identifies $L_{\rm reg}(\Delta, j+1) = L(\Delta, j+1)$ with the standard
critical $L$-value.  \eqref{eq:Delta-main} then collapses to the
classical Eichler--Shimura formula
$\omega^{\eps_j}(\Delta) = (2\pi)^{10-j}(-i)^j\binom{10}{j}\,j!\,L(\Delta, j+1)/c_S^j$.
\end{lemma}

\begin{proof}
$M_j(\Delta, \tau_0) \to 0$ exponentially because the $T$-path integrand
$\Delta(\tau)\tau^j$ decays as $|q|^1 = e^{-2\pi\im(\tau)}$ on the
segment.  Hence $P \to 0$ by Proposition~\ref{prop:Bernoulli-P}, the
$T$-segment contribution to \eqref{eq:Delta-main} drops to zero, and
the $S$-segment $[-1/\tau_0, \tau_0]$ stretches to $[0, i\infty]$.  The
Mellin identification then follows from the Hecke--Weil integral
representation of $L(\Delta, s)$.
\end{proof}

For $\Dhat$ the $\tau_0 \to i\infty$ limit is unavailable --- both
integrals individually diverge as the $q^{-1}$ Fourier mode amplifies
each $|q_0|^{-1} = e^{2\pi\im(\tau_0)} \to \infty$.  But for finite
$\tau_0 \in \HH$ both integrals are finite and \eqref{eq:Dhat-main}
yields $\eta^\pm(\Dhat)$.  The basepoint is the regulator.

\begin{remark}[Exponentially convergent $q$-series expansion]\label{rem:bfk-series}
Substituting the Fourier expansion
$f(\tau) = \sum_{n \ge -N_f}a_f(n)\,e^{2\pi i n\tau}$ into
\eqref{eq:Delta-main} (resp.\ \eqref{eq:Dhat-main}) and integrating
term-by-term via the elementary identity
\(
  \int e^{2\pi i n \tau}\tau^j\,d\tau =
  \sum_{k=0}^j \tfrac{(-1)^k j!}{(j-k)!(2\pi i n)^{k+1}}\,e^{2\pi i n\tau}\,\tau^{j-k}
\)
gives an exponentially convergent series for the periods:
\begin{equation}\label{eq:bfk-series}
  \omega^{\eps_j}(f) \;=\; \frac{(2\pi i)^{11}}{c_S^j}\sum_{n \ne 0}
                            a_f(n)\,\mathcal{K}^j(n,\tau_0)
  \;+\; (\text{constant-term piece, vanishing for } f \in S_k^!),
\end{equation}
where $\mathcal{K}^j(n, \tau_0)$ is an explicit elementary function of
$n$ and $\tau_0$ built from finite polynomial combinations of
$q_0^n = e^{2\pi i n\tau_0}$, $\widetilde q_0^{\,n} = e^{-2\pi i n/\tau_0}$,
$\tau_0^a$, $(\tau_0-1)^a$, $(-1/\tau_0)^a$, and inverse powers
$1/(2\pi i n)^{b+1}$.  Since $|q_0|, |\widetilde q_0| < 1$ for every
$\tau_0 \in \HH$, the series converges super-exponentially in $|n|$.

For $f = \Delta$, only $n \ge 1$ contributes with $a_f(n) = \tau(n)$,
yielding an exponentially convergent series for the classical
$\omega^\pm(\Delta)$ --- the explicit linear functional of
$\{\tau(n)\}_{n \ge 1}$ promised in the introduction.  For $f = \Dhat$,
the $n = -1$ term contributes one extra finite piece
$q_0^{-1}\mathcal{K}^j(-1, \tau_0)$; the $n \ge 2$ terms decay
super-exponentially in $n$ despite the
$a_{\Dhat}(n) \sim e^{4\pi\sqrt n}$ coefficient growth.  Comparing the
$\Delta$-version of \eqref{eq:bfk-series} to BFK's incomplete-gamma
expansion of $L(\Delta, s)$ \cite{BFK2011} provides a concrete sanity
check; writing the same expansion for $\eta^\pm(\Dhat)$ produces the
natural DR-analog of BFK's regularised $L$-value series, valid uniformly
across $S_{12}^!$.  Detailed term-by-term comparison of these series is
deferred to a follow-up note.
\end{remark}

\section{Cohomology and basepoint independence}\label{sec:cohomology}

\begin{theorem}[Cohomological well-definition]\label{thm:cohomology}
The class $[C^f] \in H^1(\Gamma; V_{10}\otimes\CC)$ is independent of
$\tau_0$.  For $f \in S_{12}^!$, this class is cuspidal, and Brown's
decomposition \eqref{eq:Brown-decomp} reads off
cohomologically-meaningful coordinates: $\omega^+(f) = $ coordinate of
$[C^f]$ in $[P^+]$, $\omega^-(f) = $ coordinate in $[P^-]$.
\end{theorem}

\begin{proof}
For $\tau_0, \tau_1 \in \HH$, define the basepoint-shift potential
\begin{equation}\label{eq:bp-h}
  h(X, Y;\tau_0, \tau_1) \;:=\; (2\pi i)^{11}\!\int_{\tau_0}^{\tau_1}\!
                                f(\tau)\,(X - \tau Y)^{10}\,d\tau \in V_{10}\otimes\CC.
\end{equation}
The same path-additivity / modularity calculation as in
Lemma~\ref{lem:DR-cocycle}, applied to the basepoint-shifted integrals,
gives
\begin{equation}\label{eq:bp-cocycle}
  C^f_\gamma(\tau_1) \;-\; C^f_\gamma(\tau_0) \;=\;
    -\bigl[\,h(\tau_0, \tau_1)|_\gamma \;-\; h(\tau_0, \tau_1)\,\bigr],
\end{equation}
manifestly a coboundary in $V_{10}\otimes\CC$.  Hence
$[C^f_\bullet]$ is independent of $\tau_0$.  Cuspidality of the class
for $f \in S_{12}^!$ follows from Proposition~\ref{prop:P-exists}.  The
basis-coordinate identification is \eqref{eq:Brown-decomp}.
\end{proof}

\begin{corollary}[$\tau_0$- and $j$-independence in Theorem~\ref{thm:main}]\label{cor:tau0-j-indep}
The right side of \eqref{eq:main} is independent of $\tau_0 \in \HH$
(Theorem~\ref{thm:cohomology}) and of the choice of interior
$j \in \{1, \ldots, 9\}$ within its parity sector (since all
interior-$j$ extractions read off the same basis coordinate of $[C^f]$
in \eqref{eq:Brown-decomp}).  Together with
Theorem~\ref{thm:kernels}, this completes the proof of
Theorem~\ref{thm:main}.  Numerically, $|\Delta\omega^\pm| < 3\times 10^{-38}$
and $|\Delta\eta^\pm| < 1.3\times 10^{-32}$ between
$\tau_0 = 0.3 + 1.2i$ and $\tau_0 = -0.2 + 0.9i$ \cite{Notebook}.
\end{corollary}

\begin{remark}[$C^{\prime f}_S$ in Zagier's basis]\label{rem:Zagier-decomp}
Using \eqref{eq:Brown-decomp} together with $P^+_S = -\tfrac{36}{691}P_0 + P_1$,
$P^-_S = P_2$:
\[
  C^{\prime f}_S \;=\; a_0(f)\,P_0 \;+\; \omega^+(f)\,P_1 \;+\; \omega^-(f)\,P_2,
\]
with $a_0(f) = c(f) - \tfrac{36}{691}\omega^+(f) = [C^{\prime f}_S]_{X^{10}}$
the convention-dependent Eisenstein coordinate.  At
$\tau_0 = 0.3 + 1.2i$ \cite{NumericScript}:
\begin{align*}
  C^{\prime \Delta}_S
   &\;=\; (2.752\!\times\!10^5 - 6.156\!\times\!10^5\,i)\,P_0
       + (-68\,916\,772.80959519\ldots)\,P_1
       + (-5\,585\,015.379310402\ldots\,i)\,P_2,\\
  C^{\prime \Dhat}_S
   &\;=\; (-1.532\!\times\!10^{11} + 6.163\!\times\!10^{11}\,i)\,P_0
       + (127\,202\,100\,647.1771\ldots)\,P_1
       + (10\,276\,732\,343.64913\ldots\,i)\,P_2.
\end{align*}
The $P_1, P_2$ coordinates are precisely $\omega^\pm(\Delta)$ and
$\eta^\pm(\Dhat)$.
\end{remark}

\section{Comparison with the BFK regularisation}\label{sec:bfk}

Bringmann--Fricke--Kent \cite{BFK2011} regularise the divergent integral
$\int_0^{i\infty}\!f(\tau)\tau^{s-1}\,d\tau$ for weakly holomorphic
$f \in S_k^!$ by inserting a non-holomorphic exponential factor
$e^{2\pi\nu\im(\tau)}$ and analytically continuing in $\nu$.  This
produces finite regularised $L$-values $L^*_{\rm BFK}(f, s)$,
$s = 1, \ldots, k-1$, assembled into a polynomial
$r_{\rm BFK}(f) \in V_{k-2}\otimes\CC$ via Zagier's formula.

$r_{\rm BFK}$ is finite and well-defined, but it is a \emph{different}
polynomial from the DR $C^{\prime f}_S$ when $f$ has a cusp pole.
$r_{\rm BFK}$ satisfies the involution $r_{\rm BFK}|(1 + S) = 0$ (a
direct consequence of BFK's functional equation $L^*(s) = i^k L^*(k - s)$),
but \emph{fails} the three-term cocycle relation
$r_{\rm BFK}|(1 + U + U^2) \ne 0$ ($U := ST$) by approximately
$22\times$ the typical coefficient magnitude in the numerics for
$\Dhat$ \cite{Notebook}.  Hence $r_{\rm BFK}(\Dhat) \notin W_{12}$,
whereas $C^{\prime\Dhat}_S \in W_{12}$ by construction
(Proposition~\ref{prop:P-exists}, Theorem~\ref{thm:cohomology}).

The discrepancy is structural: BFK's single-parameter regulator
$e^{iu\tau}$ has $S$-symmetry under $\tau \to -1/\tau$ (which exchanges
the two cusps $\{0, i\infty\}$ and maps the regulator at one cusp to
the regulator at the other), but no analogous symmetry for the
$U$-orbit $\{0, 1, i\infty\}$.  The $\Gamma$-cocycle relation forced by
$(ST)^3 = -I$ on genuine cocycles therefore survives the BFK
regularisation only in part.  The DR cocycle has no such asymmetry:
both contours $\gamma^{-1}\tau_0 \to \tau_0$ stay strictly inside $\HH$,
no regulator parameter is needed, and the cocycle relation holds at the
level of $\Gamma$-cocycles, not just at the level of regularised
values.  This is the structural reason Theorem~\ref{thm:main} delivers
an element of $W_{12}$ uniformly across $S_{12} \subset S_{12}^!$.

\section{Outlook: meromorphic forms, and a gauge-invariance analogy}\label{sec:outlook}

\paragraph{Extension to interior poles.}
The construction of Theorem~\ref{thm:main} uses contours that stay
strictly inside $\HH$.  For a meromorphic modular form $f$ with poles
in the interior of $\HH$ (e.g.\ $1/E_6$, with simple poles at $\tau = i$
and its $\Gamma$-orbit), the cocycle \eqref{eq:DR-def} requires contours
that avoid those poles, and shifting the basepoint moves poles across
the contour, introducing residue contributions.

For a simple pole of $f$ at $\tau_p \in \HH$ with residue
$\mathrm{Res}(f, \tau_p)$, the residue contribution to the cocycle is
\[
  2\pi i\,(2\pi i)^{11}\,\mathrm{Res}(f, \tau_p)\,(X - \tau_p Y)^{10} \;\in\; V_{10}\otimes\CC.
\]
The multiplier polynomial is $(X - \tau_p Y)^{10}$ evaluated at the
pole.  For a pole of order $m$, the contribution is a derivative
expression in $(X - \tau_p Y)^{10-i}$ for $i = 0, \ldots, m-1$.  Each
such residue polynomial is a coboundary in $V_{10}\otimes\CC$ (in the
sense that the basepoint-shift identity \eqref{eq:bp-cocycle} acquires
residue terms when poles cross), so the cohomology class $[C^f]$
remains well-defined, and the period extraction
Theorem~\ref{thm:main} extends naturally --- with the kernel formula
\eqref{eq:main} acquiring a discrete sum over residue contributions
$\sum_p (X - \tau_p Y)^j$-weights.  Working out the precise structure
of this extension is the natural successor problem to the present note.

\paragraph{A gauge-invariance analogy, and an open question.}
The cocycle $C^f$ carries more data than the cohomologically-meaningful
periods $\omega^\pm(f)$: an entire coboundary direction $\delta P$ must
be subtracted, and the choice of representative $P$ (parametrised by
the unconstrained $p_{10} \in \CC$) is irrelevant for the final answer.
This carries strong echoes of off-shell gauge data in Feynman amplitudes
with massless gauge bosons --- longitudinal polarisations, Faddeev--Popov
ghosts --- which appear in the intermediate bookkeeping but cancel in
physical (gauge-invariant) observables.  In QFT, the on-shell S-matrix
programme circumvents the gauge problem by building amplitudes directly
from gauge-invariant on-shell building blocks (helicity variables, BCFW
recursion, unitarity cuts \cite{ElvangHuang, McGadyHuang2013,
McGadyHuangChen2014}), so the gauge-redundant intermediate data never
enters: one is forced into ``fix the gauge'' only because one has made
the upstream choice to work off-shell.

The analog question for periods: does there exist a ``manifestly
cohomological'' construction that produces $\omega^\pm(f)$ directly from
the Fourier coefficients of $f$, sidestepping the intermediate cocycle
$C^f$ and coboundary $\delta P$ scaffolding entirely?  The exponentially
convergent $q$-coefficient series of Remark~\ref{rem:bfk-series} is a
partial step --- input $\{a_f(n)\}$, output $\omega^\pm$, no $C^f$ or
$P$ visible in the final formula --- but the kernels appearing in that
series were themselves derived via $\delta_T^{-1}$, so the cocycle
structure remains implicit.  A formulation in which the kernels are
constructed without going through the cocycle / coboundary intermediate
stage --- i.e.\ the modular-form analog of writing down on-shell
amplitudes directly via BCFW or unitarity rather than gauge-fixing an
off-shell Lagrangian --- would seem to require an independent
characterisation of the cuspidal-cohomology coordinates of
$H^1_{\rm cusp}(\Gamma; V_{n})$.  We flag this as an open question;
what the modular-form analog of the on-shell S-matrix programme might
be is, in our view, a worthy direction.

\begin{thebibliography}{99}

\bibitem{Brown2017}
F.\ Brown,
\emph{A class of non-holomorphic modular forms III},
arXiv:1710.07912 (2017).

\bibitem{DR2017}
N.\ Diamantis and L.\ Rolen,
\emph{Period polynomials, derivatives of $L$-functions, and zeros of polynomials},
arXiv:1707.04814 (2017).

\bibitem{BFK2011}
K.\ Bringmann, K.-H.\ Fricke, and Z.\ Kent,
\emph{Special $L$-values and periods of weakly holomorphic modular forms},
Ramanujan J.\ \textbf{24} (2011), 119--139.

\bibitem{KZ1984}
W.\ Kohnen and D.\ Zagier,
\emph{Modular forms with rational periods},
in \emph{Modular Forms} (R.A.\ Rankin, ed.), Ellis Horwood, 1984, pp.\ 197--249.

\bibitem{ElvangHuang}
H.\ Elvang and Y.-t.\ Huang,
\emph{Scattering Amplitudes in Gauge Theory and Gravity},
Cambridge University Press, 2015.

\bibitem{McGadyHuang2013}
Y.-t.\ Huang and D.\ McGady,
\emph{Consistency Conditions for Gauge Theory S Matrices from Requirements of Generalized Unitarity},
Phys.\ Rev.\ Lett.\ \textbf{112} (2014) 241601,
arXiv:1307.4065 [hep-th].

\bibitem{McGadyHuangChen2014}
W.-M.\ Chen, Y.-t.\ Huang, and D.\ A.\ McGady,
\emph{Anomalies without an action},
arXiv:1402.7062 [hep-th] (2014).

\bibitem{Notebook}
\texttt{period\_extraction\_DR\_general.ipynb} (executed Jupyter notebook).
Companion source file; implements
\texttt{extract\_periods(coeffs\_dict, tau0, k)}: the single
$q$-series-agnostic function reproducing the $\omega^\pm$, $\eta^\pm$
values of Corollaries~\ref{cor:Delta},~\ref{cor:Dhat} and the
basepoint-independence check of Corollary~\ref{cor:tau0-j-indep}, at
$45$-digit \texttt{mpmath} precision.

\bibitem{VerifyBernoulli}
\texttt{verify\_bernoulli\_formula.py}, companion script.  Confirms
agreement of the closed-form Bernoulli operator \eqref{eq:Bernoulli-P}
with the upper-triangular recursion of Proposition~\ref{prop:P-exists}
on every $p_m$, $m = 0, \ldots, 9$, to relative residual $\sim 10^{-42}$
at $45$-digit working precision.

\bibitem{VerifyL}
\texttt{verify\_omega\_plus\_formula.py}, companion script.  Evaluates
each of the nine equivalent forms of $\omega^\pm(\Delta)$ in
\eqref{eq:Delta-main} by computing $L(\Delta, s)$ via the BFK
incomplete-gamma sum, confirming agreement with Brown's reported values
to relative error $< 6\times 10^{-42}$.  Source of the
$L_{\rm reg}(\Delta, s)$ table in \S\ref{sec:specialisation}.

\bibitem{KernelScript}
\texttt{compute\_period\_kernels.py}, companion script.  Computes the
nine universal $\QQ$-polynomial kernel pairs $(k_S^j, k_T^j)$ of
Table~\ref{tab:kernels} exactly via sympy rationals, by applying the
$\delta_T^{-1}$ Bernoulli operator monomial-by-monomial as in
\eqref{eq:beta-def}.

\bibitem{NumericScript}
\texttt{compute\_numerical\_polynomials.py}, companion script.  Computes
the explicit modified cocycle values $C^{\prime f}_S$ of
Remark~\ref{rem:Zagier-decomp} at $\tau_0 = 0.3 + 1.2i$ for
$f \in \{\Delta, \Dhat\}$ at $30$-digit \texttt{mpmath} precision.
Verifies period extraction from all nine interior monomials to lockstep
agreement.

\bibitem{Companion}
Long-form companion document \texttt{analytic\_periods\_weight12.tex}
/ \texttt{.pdf}.  Same mathematical content as the present note with
extensive pedagogical discussion of cocycles, coboundaries, the
Bernoulli-operator derivation, and worked numerical polynomials.  The
present note is the math-paper-style restatement with \eqref{eq:main}
as the front matter.

\end{thebibliography}

\end{document}
